% GENERATED FILE. Edit canonical lessons, then run npm run generate:books.
% canonical-source-hash: fnv1a64:8ae76da66978e5bc
% canonical-lessons: RU-C412-peshekhod, RU-C412-velosipedist, RU-C412-ustupat, RU-C412-propuskat, RU-C412-razmetka

\chapter{Pedestrians and cyclists}
\label{ch:ru-pedestrians-and-cyclists}
\hlchaptermodality{412}

\begin{chapteropening}
\textbf{By the end of this chapter:} \emph{I can talk about sharing the road: pedestrians, cyclists, giving way, letting someone through and road markings.}

Complete the last lesson of chapter 412: I can talk about sharing the road: pedestrians, cyclists, giving way, letting someone through and road markings.
\end{chapteropening}

\section[\texorpdfstring{\ru{пешеход} (peshekhód)}{пешеход (peshekhód)}]{\ru{пешеход} (peshekhód) — a pedestrian}

\label{lesson:RU-C412-peshekhod}



\begin{quote}
Before the new one: say the Russian for to drop off, then the Russian for to fasten one's seat belt.
\end{quote}

\subsection*{You'll want to know: \ru{пешеход}}
\textbf{\ru{пешеход}} (\emph{peshekhód}) — \textquotedblleft{}a pedestrian\textquotedblright{}.

\begin{quote}
\ru{Водитель} \ru{остановился}, \ru{и} \ru{пешеход} \ru{перешёл} \ru{улицу}.
\end{quote}

\begin{quote}
\emph{Vodítel ostanovílsya, i peshekhód pereshyól úlitsu.}
\end{quote}

\begin{quote}
The driver stopped, and the pedestrian crossed the street.
\end{quote}

Masculine; built on \emph{\ru{пешком}} — on foot.

\begin{cousinweb}
\emph{\ru{пешеход}} is built from \emph{\ru{пеш}-} 'on foot', as in \emph{\ru{пешком}}, and \emph{\ru{ход}} 'a going', which is inherited from Proto-Slavic \emph{xodŭ} 'motion' (EtymDB).
\end{cousinweb}

\begin{grammarlens}[title={What you've built}]
One more word for telling, explaining and describing at B1.
\end{grammarlens}

\subsection*{Your turn}
\begin{itemize}
\raggedright
  \item \emph{Say it:} \emph{peshekhód}
  \item \emph{Say it:} \emph{peshekhód}, once more
  \item \emph{Say it:} a sentence of your own with \emph{peshekhód}
  \item \emph{Recall:} say the Russian for to drop off, then the Russian for to fasten one's seat belt, then say \emph{peshekhód} again
\end{itemize}

\begin{tcolorbox}[breakable,colback=teal!4,colframe=teal!35!black,title={Before you move on}]
What does \ru{пешеход} mean? (\textquotedblleft{}A pedestrian\textquotedblright{}.) Say it once more.
\end{tcolorbox}
\section[\texorpdfstring{\ru{велосипедист} (velosipedíst)}{велосипедист (velosipedíst)}]{\ru{велосипедист} (velosipedíst) — a cyclist}

\label{lesson:RU-C412-velosipedist}



\begin{quote}
Before the new one: say the Russian for to fasten one's seat belt, then the Russian for a pedestrian.
\end{quote}

\subsection*{You'll want to know: \ru{велосипедист}}
\textbf{\ru{велосипедист}} (\emph{velosipedíst}) — \textquotedblleft{}a cyclist\textquotedblright{}.

\begin{quote}
\ru{В} \ru{парке} \ru{много} \ru{велосипедистов}.
\end{quote}

\begin{quote}
\emph{V párke mnógo velosipedístov.}
\end{quote}

\begin{quote}
There are a lot of cyclists in the park.
\end{quote}

Feminine \emph{\ru{велосипедистка}}.

\begin{cousinweb}
\emph{\ru{велосипедист}} is built on \emph{\ru{велосипед}} 'a bicycle' with the ending \emph{-\ru{ист}}. \emph{\ru{Велосипед}} was borrowed from French \emph{vélocipède}, from Latin \emph{vēlōx} 'swift' and \emph{pēs} 'foot' (EtymDB).
\end{cousinweb}

\begin{grammarlens}[title={What you've built}]
One more word for telling, explaining and describing at B1.
\end{grammarlens}

\subsection*{Your turn}
\begin{itemize}
\raggedright
  \item \emph{Say it:} \emph{velosipedíst}
  \item \emph{Say it:} \emph{velosipedíst}, once more
  \item \emph{Say it:} a sentence of your own with \emph{velosipedíst}
  \item \emph{Recall:} say the Russian for to fasten one's seat belt, then the Russian for a pedestrian, then say \emph{velosipedíst} again
\end{itemize}

\begin{tcolorbox}[breakable,colback=teal!4,colframe=teal!35!black,title={Before you move on}]
What does \ru{велосипедист} mean? (\textquotedblleft{}A cyclist\textquotedblright{}.) Say it once more.
\end{tcolorbox}
\section[\texorpdfstring{\ru{уступать} (ustupát)}{уступать (ustupát)}]{\ru{уступать} (ustupát) — to give way; to give up (a seat)}

\label{lesson:RU-C412-ustupat}



\begin{quote}
Before the new one: say the Russian for a pedestrian, then the Russian for a cyclist.
\end{quote}

\subsection*{You'll want to know: \ru{уступать}}
\textbf{\ru{уступать}} (\emph{ustupát}) — \textquotedblleft{}to give way; to give up (a seat)\textquotedblright{}.

\begin{quote}
\ru{В} \ru{автобусе} \ru{молодёжь} \ru{уступает} \ru{место} \ru{пожилым}.
\end{quote}

\begin{quote}
\emph{V avtóbuse molodyózh ustupáyet mésto pozhilým.}
\end{quote}

\begin{quote}
On the bus young people give up their seats to the elderly.
\end{quote}

Imperfective; perfective \emph{\ru{уступить}} (\emph{\ru{уступлю}, \ru{уступишь}}). Takes the dative of the person: \emph{\ru{уступить} \ru{дорогу} \ru{пешеходу}}.

\begin{cousinweb}
\emph{\ru{уступать}} is built from \emph{\ru{у}-} 'away' and \emph{\ru{ступать}} 'to step': to step aside for someone. EtymDB lists \emph{\ru{ступать}} without a deeper source.
\end{cousinweb}

\begin{grammarlens}[title={What you've built}]
One more word for telling, explaining and describing at B1.
\end{grammarlens}

\subsection*{Your turn}
\begin{itemize}
\raggedright
  \item \emph{Say it:} \emph{ustupát}
  \item \emph{Say it:} \emph{ustupát}, once more
  \item \emph{Say it:} a sentence of your own with \emph{ustupát}
  \item \emph{Recall:} say the Russian for a pedestrian, then the Russian for a cyclist, then say \emph{ustupát} again
\end{itemize}

\begin{tcolorbox}[breakable,colback=teal!4,colframe=teal!35!black,title={Before you move on}]
What does \ru{уступать} mean? (\textquotedblleft{}To give way; to give up (a seat)\textquotedblright{}.) Say it once more.
\end{tcolorbox}
\section[\texorpdfstring{\ru{пропускать} (propuskát)}{пропускать (propuskát)}]{\ru{пропускать} (propuskát) — to let through; to miss (a bus, a lesson)}

\label{lesson:RU-C412-propuskat}



\begin{quote}
Before the new one: say the Russian for a cyclist, then the Russian for to give way.
\end{quote}

\subsection*{You'll want to know: \ru{пропускать}}
\textbf{\ru{пропускать}} (\emph{propuskát}) — \textquotedblleft{}to let through; to miss (a bus, a lesson)\textquotedblright{}.

\begin{quote}
\ru{Водитель} \ru{останавливается} \ru{и} \ru{пропускает} \ru{велосипедиста}.
\end{quote}

\begin{quote}
\emph{Vodítel ostanávlivayetsya i propuskáyet velosipedísta.}
\end{quote}

\begin{quote}
The driver stops and lets the cyclist through.
\end{quote}

Imperfective; the perfective is \emph{\ru{пропустить}} (\emph{\ru{пропущу}, \ru{пропустишь}}). \emph{\ru{Пропустить} \ru{автобус}} — to miss the bus.

\begin{cousinweb}
\emph{\ru{пропускать}} is built from \emph{\ru{про}-} 'through' and \emph{\ru{пускать}} 'to let go, to allow' (EtymDB lists \emph{\ru{пускать}} with that sense).
\end{cousinweb}

\begin{grammarlens}[title={What you've built}]
One more word for telling, explaining and describing at B1.
\end{grammarlens}

\subsection*{Your turn}
\begin{itemize}
\raggedright
  \item \emph{Say it:} \emph{propuskát}
  \item \emph{Say it:} \emph{propuskát}, once more
  \item \emph{Say it:} a sentence of your own with \emph{propuskát}
  \item \emph{Recall:} say the Russian for a cyclist, then the Russian for to give way, then say \emph{propuskát} again
\end{itemize}

\begin{tcolorbox}[breakable,colback=teal!4,colframe=teal!35!black,title={Before you move on}]
What does \ru{пропускать} mean? (\textquotedblleft{}To let through; to miss (a bus, a lesson)\textquotedblright{}.) Say it once more.
\end{tcolorbox}
\section[\texorpdfstring{\ru{разметка} (razmétka)}{разметка (razmétka)}]{\ru{разметка} (razmétka) — road markings}

\label{lesson:RU-C412-razmetka}



\begin{quote}
Before the new one: say the Russian for to give way, then the Russian for to let through.
\end{quote}

\subsection*{You'll want to know: \ru{разметка}}
\textbf{\ru{разметка}} (\emph{razmétka}) — \textquotedblleft{}road markings\textquotedblright{}.

\begin{quote}
\ru{Зимой} \ru{разметки} \ru{на} \ru{дороге} \ru{почти} \ru{не} \ru{видно}.
\end{quote}

\begin{quote}
\emph{Zimóy razmétki na doróge pochtí ne vídno.}
\end{quote}

\begin{quote}
In winter the road markings are hard to make out.
\end{quote}

Feminine and singular for all the lines: \emph{\ru{дорожная} \ru{разметка}}.

\begin{cousinweb}
\emph{\ru{разметка}} is built on \emph{\ru{разметить}} 'to mark out', from \emph{\ru{раз}-} and \emph{\ru{метить}} 'to mark', the root also found in \emph{\ru{замечать}} 'to notice'. EtymDB does not list the word.
\end{cousinweb}

\begin{grammarlens}[title={What you've built}]
One more word for telling, explaining and describing at B1.
\end{grammarlens}

\subsection*{Your turn}
\begin{itemize}
\raggedright
  \item \emph{Say it:} \emph{razmétka}
  \item \emph{Say it:} \emph{razmétka}, once more
  \item \emph{Say it:} a sentence of your own with \emph{razmétka}
  \item \emph{Recall:} say the Russian for to give way, then the Russian for to let through, then say \emph{razmétka} again
\end{itemize}

\begin{tcolorbox}[breakable,colback=teal!4,colframe=teal!35!black,title={Before you move on}]
What does \ru{разметка} mean? (\textquotedblleft{}Road markings\textquotedblright{}.) Say it once more.
\end{tcolorbox}
